\chapter{Contoh Tugas, Portofolio, dan Rubrik Refleksi}
\label{app:tugas}

Lampiran ini memuat contoh lembar tugas, template laporan, dan rubrik
refleksi yang dapat langsung dipakai atau disesuaikan.

\section{Lembar Tugas Praktikum}
\label{app:tugas-praktikum}

\begin{table}[H]
\centering
\caption{Format lembar tugas praktikum}
\label{tab:app-tugas-format}
\begin{tabularx}{\textwidth}{@{}lX@{}}
\toprule
\textbf{Bagian} & \textbf{Isi} \\
\midrule
Identitas & Nama, NIM, kelas, tanggal \\
Tujuan & Sub-CPMK yang diukur \\
Deskripsi & Konteks dan masalah yang diselesaikan \\
Keluaran & Berkas, laporan, demo \\
Kriteria & Rubrik dan bobot \\
Batas waktu & Tanggal dan waktu pengumpulan \\
Integritas & Pernyataan orisinalitas \\
\bottomrule
\end{tabularx}
\end{table}

\section{Template Laporan Proyek}
\label{app:tugas-laporan}

Laporan proyek memuat:
\begin{enumerate}[label=\arabic*.,leftmargin=*]
  \item \textbf{Pendahuluan}: tujuan, ruang lingkup, dan batasan.
  \item \textbf{Spesifikasi bahasa}: token, grammar, dan aturan
        semantik.
  \item \textbf{Arsitektur}: modul, alur data, dan keputusan desain.
  \item \textbf{Implementasi}: lexer, parser, AST, type checker,
        validator.
  \item \textbf{Pengujian}: kasus uji, hasil, dan cakupan.
  \item \textbf{Keterbatasan}: fitur yang belum didukung.
  \item \textbf{Refleksi}: pembelajaran dan rencana perbaikan.
  \item \textbf{Lampiran}: kode, keluaran, dan tautan repositori.
\end{enumerate}

\section{Template Jurnal Refleksi}
\label{app:tugas-jurnal}

\begin{table}[H]
\centering
\caption{Template jurnal mingguan}
\label{tab:app-jurnal}
\begin{tabularx}{\textwidth}{@{}lX@{}}
\toprule
\textbf{Pertanyaan} & \textbf{Jawaban} \\
\midrule
Apa yang dipelajari? & \\
Apa yang berhasil? & \\
Apa yang sulit? & \\
Apa yang akan dilakukan? & \\
Bukti/tautan & \\
\bottomrule
\end{tabularx}
\end{table}

\section{Rubrik Refleksi}
\label{app:tugas-refleksi}

\begin{table}[H]
\centering
\caption{Rubrik penilaian refleksi}
\label{tab:app-rubrik-refleksi}
\begin{tabularx}{\textwidth}{@{}lXr@{}}
\toprule
\textbf{Aspek} & \textbf{Indikator} & \textbf{Bobot} \\
\midrule
Kedalaman & Menjelaskan penyebab, bukan hanya kejadian & 30\% \\
Keterkaitan & Menghubungkan pengalaman dengan konsep & 25\% \\
Kejujuran & Mengakui kesulitan dan kesalahan & 20\% \\
Rencana & Tindakan konkret dan terukur & 25\% \\
\bottomrule
\end{tabularx}
\end{table}

\section{Contoh Entri Portofolio}
\label{app:tugas-portofolio}

\begin{casestudy}
\textbf{Judul:} Menangani dangling else pada parser Subset-C.

\textbf{Masalah:} Parser awal memasangkan \code{else} ke \code{if}
luar, sehingga program \code{if (a) if (b) x; else y;} berperilaku
salah.

\textbf{Tindakan:} Mengubah grammar dan menambahkan uji kasus.
Parser memilih shift pada token \code{else}, sesuai aturan else
terdekat.

\textbf{Bukti:} Berkas uji \file{if\_else.sc} dan keluaran AST.

\textbf{Refleksi:} Ambiguitas tidak cukup diselesaikan dengan
intuisi; perlu aturan eksplisit dan pengujian.
\end{casestudy}
